\pagebreak\section{Executive Summary}
\par \vspace{0.5cm}

The ECCO Version 4 Release 4 (V4r4) products are comprehensive global ocean and sea-ice state estimates spanning from 1992 to 2018. These datasets are dynamically and kinematically consistent reconstructions of the three-dimensional, time-evolving ocean, sea-ice, and surface atmospheric states. They include a wide range of variables such as temperature, salinity, velocity, sea level anomalies, and fluxes (e.g., heat, freshwater, and salt). The data are available at daily, monthly, and instantaneous intervals on both the high-resolution LLC90 grid and a 0.5-degree interpolated grid. The datasets adhere to modern metadata standards and are formatted in netCDF-4 (and may be available in other data format such as \hyperlink{https://en.wikipedia.org/wiki/Zarr_(data_format)}{Zarr data format}) for accessibility via NASA's Earthdata Cloud infrastructure ( \hyperlink{https://podaac.jpl.nasa.gov/}{PO.DAAC-NASA}).

The key features of ECCO V4r4 include its ability to assimilate diverse observational datasets from satellites and in situ measurements using the MIT general circulation model (MITgcm). This allows for accurate representation of global ocean dynamics and climate processes. The release also introduces advanced cloud-native services for efficient data access and processing. These products are essential for studying ocean circulation, climate variability, sea-level changes, and freshwater fluxes on a global scale.

This user guide document first provides the scope and the key overall content of the present document, then an overview of ECCO data products (Level 4) followed by detailed technical specifications of dataset filename, file structure and supporting configuration. Finaly, a full description of grid geometry, the coordinates and the data variables for both native lat-lon-cap 90 (llc90), latlon 0.5-degree and 1D datasets are provided with example of each of them.

